\labsection{3}{Аппроксимация функции с помощью двухслойной нейронной сети}

\labpart{Цель работы}
Изучить архитектуру двухслойной нейронной сети (многослойного персептрона), принципы её обучения методом
обратного распространения ошибки и исследовать точность аппроксимации нелинейной функции в зависимости
от числа нейронов скрытого слоя.

\labpart{Теоретические сведения}
Раздел~\ref{sec:mlp} теоретической части: сеть 1--$S$--1 (\ref{eq:mlp}), теорема об универсальной
аппроксимации, алгоритм обратного распространения ошибки, градиентный спуск и условие устойчивости
(\ref{eq:lr-bound}), ускоренные алгоритмы, критерии останова, недообучение и переобучение,
реализация на Python (подраздел~\ref{sec:py-mlp}).

\labpart{Постановка задачи}
Дана функция $y = f(t)$ (таблица~\ref{tab:l3-var}). Необходимо аппроксимировать её по вектору значений
$y$ для $t = t_0 : 0,1 : t_1$ (\code{np.arange(t0, t1 + 0.05, 0.1)}) двухслойной сетью 1--$S$--1 с
функцией активации $\tanh$ в скрытом слое и линейной в выходном. Сеть обучается пакетным градиентным
спуском по методу обратного распространения ошибки с параметрами \code{show = 50} (период
перерисовки), \code{lr = 0.05} (коэффициент обучения), \code{epochs = 600} (число эпох),
\code{goal = 1e-2} (целевая ошибка); в исходной постановке для MATLAB это функции \code{newff} и
\code{traingd}. Построить
графики функции и выхода сети на одной координатной плоскости и проанализировать точность
аппроксимации в зависимости от числа нейронов скрытого слоя.

\begin{table}[!htb]
\caption{Варианты заданий к лабораторной работе №3}
\label{tab:l3-var}
\centering\small
\begin{tabular}{clc|clc}
\toprule
№ & $f(t)$ & $t$ & № & $f(t)$ & $t$ \\
\midrule
1 & $(1,5t + 0,9)\sin(t + 1,1)$     & $[-5; 5]$   & 8  & $(1,5t + 0,9)\cos(t + 1,1)$     & $[-7; 7]$ \\
2 & $(1,3t + 1,9)\cos(1,1t - 1,5)$  & $[-6; 6]$   & 9  & $(0,3t + 1,1)\sin(1,8t - 0,5)$  & $[-6; 6]$ \\
3 & $(t + 1,3)\sin(0,5t + 1)$       & $[-7; 7]$   & 10 & $(0,5t + 1,9)\cos(0,7t - 1,3)$  & $[-5; 5]$ \\
4 & $(1,1t - 1,7)\cos(t + 1,5)$     & $[-9; 9]$   & 11 & $(2,5t - 0,9)\sin(1,1t + 1,7)$  & $[-6; 6]$ \\
5 & $(2,5t + 1,7)\sin(1,1t + 0,7)$  & $[-10; 10]$ & 12 & $(2,1t + 0,8)\cos(1,3t - 0,9)$  & $[-7; 7]$ \\
6 & $(0,5t - 1,4)\cos(0,5t + 1,1)$  & $[-9; 9]$   & 13 & $(0,9t + 1,5)\sin(1,1t + 1)$    & $[-8; 8]$ \\
7 & $(1,7t + 1,5)\sin(0,7t - 1,4)$  & $[-8; 8]$   & 14 & $(1,7t - 1,9)\cos(1,4t - 1,7)$  & $[-9; 9]$ \\
\bottomrule
\end{tabular}
\end{table}

\labpart{Требования к программе}
Через интерфейс задаются число нейронов скрытого слоя (ползунок от 1 до 30), коэффициент обучения
$\eta \in [0,01; 0,5]$ и целевая ошибка \code{goal}~$\in [10^{-4}; 10^{-1}]$. Точки функции и кривая
аппроксимации выводятся на одной плоскости; кривая перерисовывается каждые \code{show} эпох, синхронно
строится график ошибки по эпохам. Скорость анимации задаётся ползунком <<Задержка (мс)>>, процесс
управляется кнопками <<Старт>>, <<Пауза>>, <<Сброс>>. В Jupyter интерфейс собирается из
\code{ipywidgets} (\code{IntSlider}, \code{FloatText}, \code{Button}) и графика
\code{plotly.graph\_objects.FigureWidget}; допустимо и отдельное веб-приложение.

\labpart{Порядок выполнения работы}
\begin{steps}
\item Сформировать вектор $t = t_0 : 0,1 : t_1$ и вектор эталонных значений $y = f(t)$.
\item Создать сеть с 10 нейронами и параметрами обучения по заданию, обучить её с мультипликацией.
Проанализировать форму итоговой кривой и протокол обучения.
\item Провести серию экспериментов при $S = 1, 2, 5, 10, 20, 30$, для каждого $S$~--- не менее 5--10 запусков
с разной начальной инициализацией. Занести в таблицу медиану MSE, $R^2$ и факт достижения цели.
\item Для оценки качества между узлами обучающей сетки вычислить ошибку на контрольной сетке с шагом 0,01.
\item Дополнительно: сравнить градиентный спуск со спуском с моментом и адаптивным шагом, методом
Левенберга~--- Марквардта (\code{scipy.optimize.least\_squares}) и \code{MLPRegressor}; исследовать влияние
\code{lr} и найти границу устойчивости; добавить к значениям функции шум и показать переобучение
избыточной сети и эффект раннего останова.
\end{steps}

\labpart{Пример выполнения}
Вариант~3: $f(t) = (t + 1,3)\sin(0,5t + 1)$, $t \in [-7; 7]$, 141 точка. Сеть с $S = 10$ за 600 эпох
градиентного спуска уменьшила MSE со 145 до 0,045 ($R^2 = 0,995$); за первые 50 эпох кривая грубо
повторяет форму функции, далее медленно <<дотягивает>> экстремумы, цель $10^{-2}$ не достигнута.
Результаты серии~--- в таблице~\ref{tab:l3-series} и на рисунке~\ref{fig:l3-series}.

\begin{table}[!htb]
\caption{Серия экспериментов (медиана по 10 инициализациям): ГС~--- градиентный спуск, ГС+АШ~---
спуск с моментом и адаптивным шагом, ЛМ~--- метод Левенберга~--- Марквардта}
\label{tab:l3-series}
\centering\small
\begin{tabular}{ccccccc}
\toprule
\multirow{2}{*}{$S$} & \multirow{2}{*}{Параметров} & \multicolumn{2}{c}{ГС, 600 эпох} &
  \multicolumn{2}{c}{Эпох до цели $10^{-2}$} & Предельная MSE \\
 & & MSE & $R^2$ & ГС+АШ & ЛМ & (ЛМ) \\
\midrule
1  & 4  & 4,08  & 0,53  & ---   & ---  & 1,5 \\
2  & 7  & 2,18  & 0,75  & ---   & ---  & 0,68 \\
5  & 16 & 0,048 & 0,995 & 145   & 5    & $7 \cdot 10^{-9}$ \\
10 & 31 & 0,049 & 0,994 & 117   & 3    & $2,5 \cdot 10^{-9}$ \\
20 & 61 & 0,067 & 0,992 & 196   & 1    & $1 \cdot 10^{-9}$ \\
30 & 91 & 0,096 & 0,989 & 386   & 1    & $3 \cdot 10^{-10}$ \\
\bottomrule
\end{tabular}
\end{table}

\begin{figure}[!htb]
\centering
\includegraphics[width=0.92\textwidth]{\labfig{3}{04_series_approximations.png}}
\caption{Аппроксимация при разном числе нейронов (градиентный спуск, $\eta = 0,05$, 600 эпох)}
\label{fig:l3-series}
\end{figure}

При $S = 1, 2$~--- недообучение при любом алгоритме; цель достижима начиная с $S = 3$ (по нейрону на
изгиб функции), при $S \ge 5$ ошибка снижается до $10^{-8}$. Рябь у сетей с $S = 20, 30$~--- недоученность,
а не переобучение: ошибка на контрольной сетке равна ошибке на обучающей. Без ограничения в 600 эпох
градиентный спуск достигает цели за 2360--3960 эпох, спуск с адаптивным шагом~--- за сотни, метод
Левенберга~--- Марквардта~--- за единицы эпох. Граница устойчивости градиентного спуска: при $\eta = 0,2$ обучение сходится, при $\eta \ge 0,25$
расходится, что согласуется с оценкой $2/\lambda_{\max} \approx 0,19$ по гессиану. На данных с шумом
$\sigma = 0,7$ сеть с $S = 30$ переобучается (рисунок~\ref{fig:overfit}), лучший результат даёт $S = 3$, а
ранний останов уменьшает ошибку избыточной сети в 11 раз.

\labpart{Контрольные вопросы}
\begin{questions}
\item Какова математическая и алгоритмическая роль нелинейной функции активации в скрытом слое?
\item Опишите смысл алгоритма обратного распространения ошибки. Как вычисляются градиенты по весам
скрытого и выходного слоёв?
\item Как коэффициент обучения влияет на сходимость? Чем определяется предельное значение \code{lr}?
\item Сформулируйте теорему об универсальной аппроксимации (Цыбенко, Колмогорова). Что она говорит о
числе слоёв и чего не гарантирует?
\item Каковы визуальные признаки недообучения и переобучения на графике аппроксимации? Как отличить
переобучение от недоученности?
\item Объясните назначение параметров \code{show}, \code{goal}, \code{epochs}, \code{min\_grad}.
\item Зачем нормировать входы и цели (\code{MinMaxScaler})? Как инициализировать веса (Нгуен~---
Уидроу, Глоро)?
\end{questions}
